% concatenated from singular-constant-for-arxiv.tex
% ------------------------------------------------------------------------
% bjourdoc.tex for birkjour.cls*******************************************
% ------------------------------------------------------------------------
%%%%%%%%%%%%%%%%%%%%%%%%%%%%%%%%%%%%%%%%%%%%%%%%%%%%%%%%%%%%%%%%%%%%%%%%%%

\documentclass{birkjour}
%
%
% THEOREM Environments (Examples)-----------------------------------------
%
 \newtheorem{thm}{Theorem}[section]
 \newtheorem{cor}[thm]{Corollary}
 \newtheorem{lem}[thm]{Lemma}
 \newtheorem{prop}[thm]{Proposition}
 \theoremstyle{definition}
 \newtheorem{defn}[thm]{Definition}
 \theoremstyle{remark}
 \newtheorem{rem}[thm]{Remark}
 \newtheorem*{ex}{Example}
 \numberwithin{equation}{section}


% \usepackage{showkeys}

\def\r{\mathbb R}
\def\h{\mathbb H}
\def\v{\vec{v}}
\def\e{\mathbf e}
\def\R{\mathbb R}
\def\n{\textbf{n}}
\def\a{\textbf{a}}
\def\s{\mathbb S}
\def\M{\mathbb M}


\begin{document}

%-------------------------------------------------------------------------
% editorial commands: to be inserted by the editorial office
%
%\firstpage{1} \volume{228} \Copyrightyear{2004} \DOI{003-0001}
%
%
%\seriesextra{Just an add-on}
%\seriesextraline{This is the Concrete Title of this Book\br H.E. R and S.T.C. W, Eds.}
%
% for journals:
%
%\firstpage{1}
%\issuenumber{1}
%\Volumeandyear{1 (2004)}
%\Copyrightyear{2004}
%\DOI{003-xxxx-y}
%\Signet
%\commby{inhouse}
%\submitted{March 14, 2003}
%\received{March 16, 2000}
%\revised{June 1, 2000}
%\accepted{July 22, 2000}
%
%
%
%---------------------------------------------------------------------------
%Insert here the title, affiliations and abstract:
%


\title[Singular minimal surfaces with constant  curvature]
 {Singular minimal surfaces with constant  curvature}



%----------Author 1
 
\author{Rafael L\'opez}
\address{Department of Geometry and Topology. University of Granada. 18071  Granada, Spain}
\email{rcamino@ugr.es}
%----------classification, keywords, date
\subjclass{53A04, 53B25}
\keywords{singular minimal surface, Gauss curvature, principal curvatures, Codazzi equations}


 
\begin{abstract}
We prove that singular minimal surfaces with  constant Gauss curvature are planes, spheres and cylindrical surfaces. We also  classify all singular minimal surfaces with a constant principal curvature and  singular minimal surfaces with constant mean curvature.
\end{abstract}

%%% ----------------------------------------------------------------------
\maketitle


  

 

%%%%
\section{Introduction and statement of the results}
%%%% 

 %%%%%%%%%%%%%%%%%%%%%%%%%%%%%%%%%%%%%%%%

 Let $\vec{a}\in\r^3$ be a unit vector in Euclidean space $\r^3$ and $\alpha\in\r$. Let $\Sigma$ be an oriented surface and $\Phi\colon\Sigma\to\r^3$ be an isometric immersion of $\Sigma$ in $\r^3$. The surface $\Sigma$ is said to be an {\it $\alpha$-singular  minimal surface} if its mean curvature $H$ satisfies
 \begin{equation}\label{eq1}
 H=\alpha\frac{\langle N,\vec{a}\rangle}{\langle\Phi,\vec{a}\rangle},
 \end{equation}
 where $N$ is the unit normal vector of $\Sigma$. We are implicitly assuming that  $\Phi(\Sigma)$ is included in one of the two halfspaces determined by the vector plane $\Pi=\{p\in\r^3\colon\langle  p,\vec{a}\rangle=0\}$. Singular minimal surfaces are   critical points of the energy functional 
 $$\Sigma\mapsto\int_\Sigma \langle\Phi(p),\vec{a}\rangle^\alpha\, d\Sigma,$$
 where $d\Sigma$ is the area element of $\Sigma$. If $\alpha=0$, this energy is simply the area functional and the critical points are minimal surfaces. From now, we discard the case $\alpha=0$. The case $\alpha=1$ extends in dimension two the notion of catenary: see \cite{bht,dh,ni}. The study of the singular minimal surface equation can be viewed in \cite{bd,di1,di2,lo3,lo4}. For  geometric aspects of singular minimal surfaces, see \cite{dg,dl,lo1,lo2}.  

Among the examples of $\alpha$-singular minimal surfaces, we focus on isoparametric surfaces, that is,  planes, spheres and circular cylinders. 
\begin{enumerate}
\item Planes parallel to $\vec{a}$ are $\alpha$-singular minimal surfaces for all $\alpha\in\r$. 
\item Spheres are $\alpha$-singular minimal surfaces only if $\alpha=-2$. In such a case, the center of the sphere lies in the vector plane $\Pi$. These surfaces also represent minimal surfaces in hyperbolic space $\mathbb{H}^3$ when $\h^3$ is viewed in the  upper halfspace model   $\r^3_+=\{p\in\r^3\colon\langle p,\vec{a}\rangle>0\}$. 
\item Circular cylinders are   $\alpha$-singular minimal surfaces only if $\alpha=-1$. In such a case, the rotational axis lies in the vector plane $\Pi$. 
\end{enumerate}


Other examples of $\alpha$-singular minimal surfaces can be found in the family of  cylindrical surfaces, that is, ruled surfaces whose rulings are parallel straight-lines. If the cylindrical surface is parametrized by $\Psi(s,t)=\gamma(s)+t\vec{v}$, $|\vec{v}|=1$, and if $\Sigma$ satisfies \eqref{eq1}, then either $\Sigma$ is a plane parallel to $\vec{a}$ or $\langle \vec{v},\vec{a}\rangle=0$. This implies that the generating curve $\gamma$ lies contained in a plane parallel to $\vec{a}$. In such a case, Eq. \eqref{eq1} writes in terms of the curvature $\kappa$ of $\gamma$ as
 \begin{equation}\label{eq2}
 \kappa= \alpha \frac{\langle{\bf n},\vec{a}\rangle}{\langle\gamma,\vec{a}\rangle},
 \end{equation}
where ${\bf n}$ is the unit normal of $\gamma$.  These surfaces have zero   Gauss curvature. Cylindrical singular minimal surfaces have been classified in \cite{dl,lo1}. If $\alpha=-1$, cylindrical singular minimal surfaces are circular cylinders. In this paper, we will refer these surfaces simply by cylindrical singular minimal surfaces.


Planes parallel to $\vec{a}$,  spheres centered at $\Pi$ and the cylindrical singular minimal surfaces  have   constant Gauss curvature $K$. The purpose of this paper is to prove the converse and characterize these surfaces are the only singular minimal surfaces with constant Gauss curvature. 
   

\begin{thm} \label{t1}
Planes parallel to $\vec{a}$,  spheres centered at $\Pi$ and the cylindrical singular minimal surfaces are the only $\alpha$-singular minimal surfaces with   constant Gauss curvature. 
\end{thm}
This result will proved in Sect. \ref{s3}.  We also study those singular minimal surfaces with a constant principal curvature.  In Sect. \ref{s5}, we prove the following classification. 

\begin{thm} \label{t22} Planes parallel to $\vec{a}$, spheres centered at $\Pi$ and cylindrical singular minimal surfaces      are the only  singular minimal surfaces with a constant principal curvature.   
\end{thm}

Finally we consider  singular minimal surfaces with constant mean curvature.

\begin{thm} \label{t3}  Planes parallel to $\vec{a}$ and spheres centered at $\Pi$ are the only   singular minimal surfaces with constant mean curvature.   
\end{thm}



%%%%%%%%%
\section{Preliminaries}\label{s2}
%%%%%%%%%%

In this section we describe parametrizations of a surface by lines of curvature. This will be done in open sets of a surface free of umbilic points. Let $\Sigma$ be an orientable surface immersed in $\r^3$. Denote by  $\nabla$ and $\overline{\nabla}$ be the Levi-Civita connections of $\Sigma$ and
$\mathbb{R}^3$, respectively.  Denote by $N$ the unit normal vector field of $\Sigma$. Let $\mathfrak{X}(\Sigma)$ be the space of tangent vector fields of $\Sigma$. The  Gauss and Weingarten formulas are, respectively, 
\begin{equation}\label{gn}  
\begin{split}
    \overline{\nabla}_XY&=\nabla_XY+h(X,Y), \\
    \overline{\nabla}_XN&=-SX
    \end{split}
\end{equation}
for all $X,Y\in\mathfrak{X}(\Sigma)$, where $h\colon T\Sigma\times T\Sigma\to T^\perp\Sigma$ and  $S\colon T\Sigma\to  T\Sigma$ are the second fundamental form and the shape
operator, respectively. The  Gauss
equation and Codazzi equations are 
\begin{equation*}
\begin{split}
    \langle R(X,Y)W,Z\rangle&=\langle h(Y,W),h(X,Z)\rangle-\langle
    h(X,W),h(Y,Z)\rangle\\
    (\nabla_XS)Y&=(\nabla_YS)X,
    \end{split}
\end{equation*}
where $R$ is the curvature tensor of $\nabla$. The   mean curvature $H$ and the Gauss curvature $K$ of $\Sigma$ are defined by  
\[ H= \kappa_1+\kappa_2 \quad K=\kappa_1\kappa_2,\]
where   $\kappa_1$ and $\kappa_2$ are the principal curvatures of $\Sigma$. In an open set $\Omega\subset\Sigma$ of non-umbilic points, consider a local
orthonormal frame $\{e_1,e_2\}$ given by  coordinates of  lines of curvature. The expression of the shape operator $S$ is 
\begin{equation}\label{b0}
Se_1 =\kappa_1 e_1,\qquad Se_2 =\kappa_2 e_2.
    \end{equation}
Since $\{e_1,e_2\}$ is an orthonormal frame, if $X\in\mathfrak{X}(\Sigma)$ then
\begin{equation*}
\nabla_X e_1=\omega(X)e_2\quad  \nabla_X e_2=-\omega(X)e_1,
\end{equation*}
where $\omega(X)=\langle\nabla_X e_1,e_2\rangle=-\langle\nabla_X e_2,e_1\rangle$. Then we have 
\begin{equation}\label{eiej}
\left\{
\begin{split}
&\nabla_{e_1}e_1=\omega(e_1)e_2,\quad \nabla_{e_1}e_2=-\omega(e_1)e_1\\
&\nabla_{e_2}e_1=\omega(e_2)e_2,\quad \nabla_{e_2}e_2=-\omega(e_2)e_1.
\end{split}\right.
\end{equation}
Using the Codazzi equations, we have  
\begin{equation}\label{b1}
    \left\{
    \begin{array}{ll}
        e_{2}(\kappa_1)=(\kappa_1-\kappa_2)\omega(e_1)\\
        e_{1}(\kappa_2)=(\kappa_1-\kappa_2)\omega(e_2),
    \end{array}
    \right.
\end{equation}
where, as usually, $e_i(f)$ denotes the derivative of a smooth function $f$ along the vector $e_i$, $i=1,2$.  Notice that $\kappa_1-\kappa_2\not=0$ in $\Omega$.
Thanks to  \eqref{b1}, in the set $\Omega$, the  Gauss equation  gives the following expression of the Gauss curvature $K$,  
\begin{equation}\label{bb}
   K=\kappa_1\kappa_2=-e_1\left(\frac{e_1( \kappa_2 )}{\kappa_1-\kappa_2} \right)+e_2\left(\frac{e_2( \kappa_1 )}{\kappa_1 -\kappa_2}\right)
   -\frac{e_1( \kappa_2)^2+ e_2( \kappa_1 )^2}{(\kappa_1 -\kappa_2)^2}.
\end{equation}
 
 

 Denote by $\Phi\colon\Sigma\to\r^3$ the immersion of  a singular minimal surface. Let  decompose the  vector $\vec{a}$ in its tangent part $\vec{a}^\top$ and normal part $\vec{a}^\perp$ with respect to $\Sigma$,  
\begin{equation*}
\vec{a}^\top=\gamma e_1+\mu e_2,
\end{equation*}
for some smooth functions $\gamma$ and $\mu$. The functions $\gamma$ and $\mu$ are given by 
$$\gamma=\langle \vec{a},e_1\rangle,\quad \mu=\langle\vec{a},e_2\rangle.$$
 The normal part $\vec{a}^\perp$  is   $\langle N,\vec{a}\rangle N$, which by \eqref{eq1} is
$$\vec{a}^{\perp}=\frac{H\langle\Phi,\vec{a}\rangle}{\alpha}N.$$
\begin{lem}\label{l2}
Let $\Omega$ be a subset of non-umbilic points of an $\alpha$-singular minimal surface. Then the functions   $\gamma$ and $\mu$ satisfy the following equations
    \begin{equation}\label{d3}
    \left\{
\begin{split}
            0=&e_{1}(\gamma)-\mu\omega(e_1)-\frac{H\langle\Phi,\vec{a}\rangle}{\alpha}\kappa_1\\
            0=&e_{2}(\gamma)-\mu\omega(e_2) \\
            0=&e_{1}(\mu)+\gamma\omega(e_1)\\
           0= &e_{2}(\mu)+\gamma\omega(e_2)-\frac{H\langle\Phi,\vec{a}\rangle}{\alpha}\kappa_2\\
           0= &e_{1}(\frac{H\langle\Phi,\vec{a}\rangle}{\alpha})+\gamma\kappa_1\\
           0= &e_{2}(\frac{H\langle\Phi,\vec{a}\rangle}{\alpha})+\mu\kappa_2.
        \end{split}\right.
    \end{equation}
\end{lem}

\begin{proof} 
 If $X\in\mathfrak{X}(\Sigma)$, then   $\overline{\nabla}_X\vec{a}=0$ because $\vec{a}$ is constant in $\r^3$. Then the Gauss and
Weingarten formulas imply 
\begin{equation*}
\begin{split}
    0&=\overline{\nabla}_{X}\vec{a}=\overline{\nabla}_{X}(\vec{a}^{\top}+\vec{a}^{\perp})=\overline{\nabla}_X\vec{a}^\top+\overline{\nabla}_X\vec{a}^\perp\\
    &={\nabla}_{X}\vec{a}^{\top}+h(X,\vec{a}^{\top})+X( H\frac{\langle\Phi,\vec{a}\rangle}{\alpha})N- H\frac{\langle \Phi,\vec{a}\rangle}{\alpha} SX 
\end{split}
\end{equation*}
for all   $X\in\mathfrak{X}(\Sigma)$.
Taking the tangent and the normal part in both sides of the above expression, we obtain  
\begin{equation}\label{tn} 
    \left\{
    \begin{array}{ll}
        {\nabla}_{X}\vec{a}^{\top}-H\frac{\langle\Phi,\vec{a}\rangle}{\alpha}SX=0\\
        h(X, \vec{a}^{\top})+X(H\frac{\langle\Phi,\vec{a}\rangle}{\alpha})N=0.
    \end{array}
    \right.
\end{equation}
Equations \eqref{d3} are obtained by substituting  in the two previous identities $X$ by $e_1$ and   $e_2$  and writing in coordinates with respect to $\{e_1,e_2,N\}$. Indeed, if $X=e_1$, and using \eqref{eiej}, the first equation in \eqref{tn} is 
\begin{equation*} 
    \begin{array}{ll}
  0&=  {\nabla}_{e_1}(\gamma e_1+\mu e_2)-H\frac{\langle \Phi,\vec{a}\rangle}{\alpha}\kappa_1 e_1\\
     & =e_1(\gamma)e_1+e_1(\mu)e_2+\gamma \omega(e_1)e_2-\mu\omega(e_1)e_1-H\frac{\langle \Phi,\vec{a}\rangle}{\alpha}\kappa_1 e_1.
         \end{array}
\end{equation*}
Similarly, the second equation in \eqref{tn} is 
$$   0 = h(e_1,\gamma e_1+\mu e_2)+e_1(H\frac{\langle\Phi,\vec{a}\rangle}{\alpha})N =\left(\gamma\kappa_1+e_1\left(\frac{ H\langle\Phi,\vec{a}\rangle}{\alpha}\right)\right)N.$$
This gives the first, third and fifth equations in \eqref{d3}. The other equations of \eqref{d3} are obtained by putting $X=e_2$ in \eqref{tn}.

\end{proof}

%%%%%%%%%%
\section{Proof of Theorem \ref{t1}}\label{s3}
%%%%%%%%%%%%%%%%%%%%%%%%%%%%%%%
We know that planes parallel to $\vec{a}$,  spheres centered at $\Pi$ and   cylindrical singular minimal surfaces have  constant Gauss curvature. We prove  the converse. 

The case $K=0$ is studied separately. If $K=0$, the surface is also called developable. It is known that these surfaces can be parametrized locally as ruled surfaces \cite{hn}. On the other hand, ruled singular minimal surfaces were classified in \cite{ay}, proving that these surfaces must be cylindrical surfaces. This proves Thm. \ref{t1} when $K=0$.


Suppose now that $K\not=0$. Let $\Sigma$ be an $\alpha$-singular minimal surface with non-zero constant Gauss curvature $K$, $K=c$. If a principal curvature is constant, and because $K\not=0$, then the other principal curvature is constant. This proves that $\Sigma$ is an isoparametric surface. Since $K\not=0$, then $\Sigma$ is a sphere. This proves the result in this situation. 

Suppose now that the two principal curvatures are not constant and we will arrive to a contradiction.  We write the principal curvature $\kappa_2$ and the mean curvature   $H$ in terms of $K$ and the principal curvature  $\kappa_1$. So, we have  
\begin{equation}\label{gm2}
\kappa_2 =\frac{c}{\kappa_1},\qquad H =\frac{\kappa_1^2+c}{\kappa_1},\qquad 
\kappa_1-\kappa_2=\frac{\kappa_1^2-c}{\kappa_1}.
\end{equation}
We now write Eq. \eqref{bb}. Using that $K=c$ is constant, and taking into account \eqref{gm2}, we have  
 \begin{equation}\label{bb2}
  \begin{split}
    c\frac{(\kappa_1^2-c)^2}{\kappa_1^2}&=\frac{c(\kappa_1^2-c)}{\kappa_1^3}e_{11}(\kappa_1) +\frac{\kappa_1^2-c}{\kappa_1}e_{22}(\kappa_1)\\
   &-\frac{3c}{\kappa_1^2}e_1(\kappa_1)^2-\frac{2\kappa_1^2+c}{\kappa_1^2}e_2(\kappa_1)^2.
   \end{split}
\end{equation}
By using the expressions of $H$ and $\kappa_2$ in \eqref{gm2},  we have
\begin{equation}\label{gm22}
\begin{split}
e_1(\frac{H\langle\Phi,\vec{a}\rangle}{\alpha})&= \frac{\kappa_1^2-c}{\alpha\kappa_1^2}\langle\Phi,\vec{a}\rangle e_1(\kappa_1)+\gamma\frac{\kappa_1^2+c}{\alpha \kappa_1},\\
e_2(\frac{H\langle\Phi,\vec{a}\rangle}{\alpha})&= \frac{\kappa_1^2-c}{\alpha\kappa_1^2}\langle\Phi,\vec{a}\rangle e_2(\kappa_1)+\mu\frac{\kappa_1^2+c}{\alpha \kappa_1}.\\
\end{split}
\end{equation}
Substituting \eqref{gm22} into  the last two equations of \eqref{d3},   we obtain
\begin{equation*}
\begin{split}
 0&=\frac{\kappa_1^2-c}{\alpha\kappa_1^2}\langle\Phi,\vec{a}\rangle e_1(\kappa_1)+\left(\frac{\kappa_1^2+c}{\alpha \kappa_1} +\kappa_1\right)\gamma,\\
0 &= \frac{\kappa_1^2-c}{\alpha\kappa_1^2}\langle\Phi,\vec{a}\rangle e_2(\kappa_1)+\left(\frac{\kappa_1^2+c}{\alpha \kappa_1}+\frac{c}{\kappa_1}\right)\mu.\\
\end{split}
\end{equation*}
From these equations, we get expressions of $\gamma$ and $\mu$,
\begin{equation}\label{gm}
 \gamma=\frac{\langle\Phi,\vec{a}\rangle(c-\kappa_1^2)}{\kappa_1  (c+(1+\alpha)\kappa_1^2)}e_1(\kappa_1),\quad 
\mu=\frac{\langle\Phi,\vec{a}\rangle(c-\kappa_1^2 )}{ \kappa_1 (\kappa_1^2+(1+\alpha)c)}e_2(\kappa_1).
\end{equation}
 
Since Thm. \ref{t1} is local, we can assume that the denominators in \eqref{gm} are not zero in some open set of $\Omega$ because, otherwise, $\kappa_1$ would be a constant function. This also holds in the next computations. We are also using in \eqref{gm} that $c\not=0$. This implies that the parameter $\alpha$ is arbitrary in the denominator of $\gamma$. The hypothesis $c\not=0$ will be employed later in successive simplifications of equations that are multiplied by $c$. In the next and successive   computations, we use a symbolic software as {\sc Mathematica} to manipulate all expressions \cite{wo}. 

We differentiate  $\gamma$ and $\mu$ in \eqref{gm} with respect to $e_1$ and $e_2$ and next, we compare with the corresponding derivatives given in \eqref{d3}.  Notice that  $e_1(\langle\Phi,\vec{a}\rangle)= \gamma$ and  $e_2(\langle\Phi,\vec{a}\rangle)= \mu$. Then the derivatives of the functions \eqref{gm} are
\begin{equation}\label{d4}
\begin{split}
 e_1(\gamma)&=\frac{\langle\Phi,\vec{a}\rangle(c-\kappa_1^2)}{\kappa_1  (c+(1+\alpha)\kappa_1^2)}e_{11}(\kappa_1)+\frac{ \langle \Phi,\vec{a}\rangle(2+\alpha)(\kappa_1^2-3 c)}{(c+(1+\alpha)\kappa_1^2)^2}e_1(\kappa_1)^2\\ 
 e_1(\mu)&=\frac{\langle\Phi,\vec{a}\rangle(c-\kappa_1^2 )}{ \kappa_1 (\kappa_1^2+(1+\alpha)c)}e_{12}(\kappa_1)+\frac{\langle\Phi,\vec{a}\rangle(2+\alpha)(c+\kappa_1^2)(\kappa_1^2-(3+\alpha)c)}{ (\kappa_1^2+(1+\alpha)c)^2(c+(1+\alpha)\kappa_1^2)}e_1(\kappa_1)e_2(\kappa_1)\\
 e_2(\mu)&=\frac{\langle\Phi,\vec{a}\rangle(c-\kappa_1^2 )}{ \kappa_1 (\kappa_1^2+(1+\alpha)c)}e_{22}(\kappa_1)+\frac{\langle\Phi,\vec{a}\rangle(2\kappa_1^4-\alpha c^2-(6+\alpha)c\kappa_1^2)}{  \kappa_1^2(\kappa_1^2+(1+\alpha)c)^2}e_2(\kappa_1)^2.
 \end{split}
 \end{equation}
We work with  equations   \eqref{d3}. Using \eqref{b1} and \eqref{gm2}, we have
\begin{equation*}
\omega(e_1)=\frac{\kappa_1}{\kappa_1^2-c}e_2(\kappa_1),\quad \omega(e_2)=-\frac{c}{\kappa_1(\kappa_1^2-c)}e_1(\kappa_1).\\
\end{equation*}
Inserting this in \eqref{gm2}, we obtain
    \begin{equation}\label{d5}
    \left\{
\begin{split}
e_{1}(\gamma)&=\frac{\mu\kappa_1}{\kappa_1^2-c}e_2(\kappa_1)+\frac{\kappa_1^2+c}{\alpha }\langle\Phi,\vec{a}\rangle   \\
 e_{1}(\mu)&=-\frac{\gamma\kappa_1}{\kappa_1^2-c}e_2(\kappa_1)\\
 e_{2}(\mu)&=  \frac{c\gamma}{\kappa_1(\kappa_1^2-c)}e_1(\kappa_1)+\frac{\kappa_1^2+c}{\alpha\kappa_1^2}  c \langle\Phi,\vec{a}\rangle
      \end{split}\right.
   \end{equation}

From \eqref{d4} and \eqref{d5}, we get the expressions of $e_{11}(\kappa_1)$, $e_{12}(\kappa_1)$ and $e_{22}(\kappa_1)$. Using also the expressions \eqref{gm} of $\gamma$ and $\mu$, we have

\begin{equation}\label{d7}
\begin{split}
e_{11}(\kappa_1)&=\frac{(\alpha +2) k \left(\kappa_1^2-3 c\right)}{\left(\kappa_1^2-c\right) \left((\alpha +1) \kappa_1^2+c\right)}e_1(\kappa_1)^2+\frac{\kappa_1 \left((\alpha +1) \kappa_1^2+c\right)}{\left(\kappa_1^2-c\right) \left(\kappa_1^2+(1+\alpha)c\right)}e_2(\kappa_1)^2\\
&-\frac{\kappa_1 \left(k^2+c\right) \left((\alpha +1) \kappa_1^2+c\right)}{\alpha  \left(\kappa_1^2-c\right)},\\
e_{12}(\kappa_1)&=  \left(  \frac{2\kappa_1}{\kappa_1^2+(1+\alpha)c}+\frac{3\kappa_1}{c-\kappa_1^2} -\frac{2 c}{\kappa_1((\alpha +1) \kappa_1^2+c)}+\frac{2}{\kappa_1}\right) e_1(\kappa_1)e_2(\kappa_1),\\
e_{22}(\kappa_1)&=\frac{c \left(\kappa_1^2+(1+\alpha)c\right)}{\kappa_1 \left(\kappa_1^2-c\right) \left((\alpha +1) \kappa_1^2+c\right)}e_1(\kappa_1)^2+\frac{-2 \kappa_1^4+(\alpha +6) \kappa_1^2 c+\alpha  c^2}{\kappa_1 \left(c-\kappa_1^2\right) \left(\kappa_1^2+(1+\alpha)c\right)}e_2(\kappa_1)^2\\
&-\frac{c \left(\kappa_1^2+c\right) \left(\kappa_1^2+(1+\alpha)c\right)}{\alpha  \kappa_1 \left(\kappa_1^2-c\right)}.
\end{split}
\end{equation}
We distinguish the cases that $e_1(\kappa_1)=0$ identically and $e_1(\kappa_1)\not=0$.
\begin{enumerate}
\item Case  $e_1(\kappa_1)=0$ identically. Since $\kappa_1$ is not a constant function, then  $e_2(\kappa_1)\not=0$   in an open set of $\Omega$.     The first equation in \eqref{d7} is $e_{11}(\kappa_1)=0$. This gives 
$$\left(c+(1+\alpha)  \kappa_1^2\right) \left(\frac{e_2(\kappa_1)^2}{\kappa_1^2+(1+\alpha)c}-\frac{\kappa_1^2+c}{\alpha }\right)=0.$$
Thus
\begin{equation}\label{e22}
e_2(\kappa_1)^2=\frac{(\kappa_1^2+c)(\kappa_1^2+(1+\alpha)c)}{\alpha}.
\end{equation}
Differentiating with respect to $e_2$, 
$$2e_2(\kappa_1)e_{22}(\kappa_1)=\frac{4\kappa_1^3+2(2+\alpha)c\kappa_1}{\alpha}e_2(\kappa_1).$$
Simplifying by $e_2(\kappa_1)$, we have 
$$e_{22}(\kappa_1)=\frac{2\kappa_1^3+(2+\alpha)c\kappa_1}{\alpha}.$$
On the other hand, using \eqref{e22},  the last equation in \eqref{d7}  is
\begin{equation*}
\begin{split}
e_{22}(\kappa_1)&=\frac{-2 \kappa_1^4+(\alpha +6) \kappa_1^2 c+\alpha  c^2}{\kappa_1 \left(c-\kappa_1^2\right) \left(\kappa_1^2+(1+\alpha)c\right)}e_2(\kappa_1)^2-\frac{c \left(\kappa_1^2+c\right) \left(\kappa_1^2+(1+\alpha)c\right)}{\alpha  \kappa_1 \left(\kappa_1^2-c\right)}\\
&=\frac{\left(\kappa_1^2+c\right) \left(2 \kappa_1^4-(\alpha +7) \kappa_1^2 c-(2 \alpha +1) c^2\right)}{\alpha  \kappa_1 \left(\kappa_1^2-c\right)}
\end{split}
\end{equation*}
  Equating the two previous  expressions of $e_{22}(\kappa_1)$, we get 
  $$ (2 \alpha +5) \kappa_1^4+2 (\alpha +3) \kappa_1^2 c+(2 \alpha +1) c^2  =0.$$
   This is a polynomial on $\kappa_1$, whose coefficients are constant and not all zero.  This implies that $\kappa_1$ is constant, obtaining a contradiction.
 \item Case $e_1(\kappa_1)\not=0$.
 We substitute \eqref{d7} in the expression of $K$ in \eqref{bb2}, obtaining
\begin{equation}\label{pe1}
P_1e_1(\kappa_1)^2+Q_1e_2(\kappa_1)^2+R_1=0,
\end{equation}
where
\begin{equation*}
\begin{split}
P_1&=\frac{\alpha  \kappa_1^2+(\alpha +4) c}{(\alpha +1) \kappa_1^2+c},\\
Q_1&= \frac{(\alpha +4) \kappa_1^2+\alpha  c}{\kappa_1^2+(1+\alpha)c},\\
R_1&= \kappa_1^4+\frac{\left(\kappa_1^2+c\right)^2}{\alpha }+c^2.
\end{split}
\end{equation*}


We differentiate \eqref{pe1} with respect to $e_1$ obtaining
$$e_1(P_1)e_1(\kappa_1)^2+e_1(Q_1)e_2(\kappa_1)^2+e_1(R_1)+2P_1e_1(\kappa_1)e_{11}(\kappa_1)+2Q_1 e_2(\kappa_1)e_{12}(\kappa_1)=0.$$
Simplifying by $e_1(\kappa_1)$ because $e_1(\kappa_1)\not=0$, we obtain 
\begin{equation}\label{pe2}
P_2e_1(\kappa_1)^2+Q_2e_2(\kappa_1)^2+R_2=0,
\end{equation}
 where 
\begin{equation*}
\begin{split}
P_2&= \frac{2 (\alpha +2) \kappa_1 \left(\alpha  \kappa_1^4+(2-3 \alpha ) k^2 c-2 (\alpha +5) c^2\right)}{\left(\kappa_1^2-c\right) \left((\alpha +1) \kappa_1^2+c\right)^2},\\
Q_2&= \frac{2 (\alpha +2) \kappa_1 \left(2 (\alpha +1) \kappa_1^6+\left(\alpha ^2-3 \alpha -8\right) \kappa_1^4 c-\left(3 \alpha ^2+12 \alpha +10\right) \kappa_1^2 c^2-\alpha  (2 \alpha +3) c^3\right)}{\left(\kappa_1^2-c\right) \left(\kappa_1^2+(\alpha+1)c\right)^2 \left((\alpha +1) \kappa_1^2+c\right)} ,\\
R_2&=\frac{2 (\alpha +2) \kappa_1^5-8 (\alpha +1) \kappa_1^3 c-2 (\alpha +6) c\kappa_1^2}{\alpha  \left(\kappa_1^2-c\right)} .
\end{split}
\end{equation*}
Before the next computations, we explain how we will arrive to a contradiction.  From Eqs. \eqref{pe1} and \eqref{pe2} we get   the values of $e_1(\kappa_1)^2$ and $e_2(\kappa_1)^2$. Here we   need to discuss if the determinant of the coefficients of $e_1(\kappa_1)^2$ and $e_2(\kappa_1)^2$, namely, $P_1Q_2-P_2Q_1$, is $0$ or not. If $P_1Q_2-P_2Q_1=0$, then a suitable combination of \eqref{pe1} and \eqref{pe2} gives a linear combination of $R_1$ and $R_2$ which will be a polynomial on $\kappa_1$ with constant coefficients. This proves that $\kappa_1$ is constant, which it is a contradiction. If  $P_1Q_2-P_2Q_1\not= 0$, and with the values of $e_1(\kappa_1)^2$ and $e_2(\kappa_1)^2$, we differentiate $e_1(\kappa_1)^2$ with respect to $e_1$. This gives an identity involving $e_{11}(\kappa_1)$, $e_{12}(\kappa_1)$, $e_1(\kappa_1)^2$ and $e_2(\kappa_1)^2$. Substituting all these functions by their  expressions, we finally obtain a polynomial on $\kappa_1$ with constant coefficients, obtaining the desired   contradiction. 

 We have $P_1Q_2-P_2Q_1=(\alpha^2 -4)   \kappa_1 \left(\kappa_1^2-c\right)^3$. Notice that $\kappa_1^2-c\not=0$.
\begin{enumerate}
\item Case $\alpha^2-4=0$.   If $\alpha=-2$, the linear combinations $P_2\eqref{pe1}-P_1 \eqref{pe2}=0$ gives $8c\kappa_1=0$, which it is not possible. If $\alpha=2$, the same  combination of both equations yields $8 c\kappa_1 \left(-5 \kappa_1^4+6 \kappa_1^2 c+15 c^2\right)=0$, which implies that $\kappa_1$ is constant, obtaining a contradiction. 
  
 \item Case $\alpha^2-4\not=0$.  Solving \eqref{pe1} and \eqref{pe2}, we obtain 
\begin{equation}\label{es2}
\begin{split}
e_1(\kappa_1)^2&:=Z_1=-\frac{M_1\left((\alpha +1) \kappa_1^2+c\right)}{\alpha ^2 \left(\alpha ^2-4\right) \left(\kappa_1^2-c\right)^3},\\
e_2(\kappa_1)^2&:=Z_2=\frac{M_2 c \left(\kappa_1^2+(1+\alpha)c\right)^2}{\alpha ^2 \left(\alpha ^2-4\right) \left(\kappa_1^2-c\right)^3},
\end{split}
\end{equation}
where 
\begin{equation*}
\begin{split}
M_1&=    (\alpha +1) (\alpha^2 -4) \kappa_1^8-(\alpha  (\alpha  (4 \alpha +11)+16)+12) \kappa_1^6 c\\
&+(\alpha  (\alpha  (7 \alpha +13)+4)-12) \kappa_1^4 c^2+(\alpha  (\alpha  (4 \alpha  (\alpha +5)+39)+16)-4) \kappa_1^2 c^3\\
&+2 \alpha ^2 (\alpha +1) (\alpha +3) c^4,\\
M_2&=   -(\alpha  (5 \alpha +8)+4) \kappa_1^4+2 (\alpha  (\alpha  (2 \alpha +5)-4)-4) \kappa_1^2 c+(\alpha  (\alpha  (2 \alpha +15)+24)-4) c^2,
\end{split}
\end{equation*}
Differentiating $e_1(\kappa_1)^2$ with respect to $e_1$, we obtain 
$$2 e_1(\kappa_1)e_{11}(\kappa_1)=e_1(\kappa_1)\frac{d}{d\kappa_1}\left(Z_1\right).$$
  Simplifying by $e_1(\kappa_1)$ and using \eqref{d7} and \eqref{es2}, we obtain 
$$ \kappa_1^2 \left(\kappa_1^2+(1+\alpha)c\right) \left((\alpha +4) \kappa_1^2+\alpha  c\right)=0.$$
This implies that $\kappa_1$ is constant, which it is a contradiction. This completes the proof of Thm. \ref{t1}.
\end{enumerate}
 \end{enumerate}

  
 %%%%%%%%
\section{Proof of Theorem \ref{t22}}\label{s5}
%%%%%%

 
It remains to prove the converse. Let $\Sigma$ be an $\alpha$-singular minimal surface with a constant principal curvature.    Without loss of generality, assume that  the principal curvature $\kappa_2$ is constant. Let $\kappa_2=c$. If $c=0$, then $K=0$ and Thm. \ref{t1} concludes that the surface is a plane or a cylindrical  singular minimal  surface. This proves Thm. \ref{t22} in this situation. 
 
 Suppose now $c\not=0$. If the principal curvature   $\kappa_1$ is constant, then the surface is isoparametric, hence it is a plane or a sphere. This proves Thm. \ref{t22} in this situation.
 
 Suppose that $\kappa_1$ is not constant and we will arrive to a contradiction. We have $K=c\kappa_1$ and $H=\kappa_1+c$. Identities \eqref{b1} gives $\omega(e_2)=0$. The expression of $K$ in \eqref{bb} is
 \begin{equation}\label{f1}
 c\kappa_1=\frac{e_{22}(\kappa_1)}{\kappa_1-c}-2\frac{e_2(\kappa_1)^2}{(\kappa_1-c)^2}. 
 \end{equation}
The last two equations of \eqref{d3} are
 \begin{equation*}
 \begin{split}
 0&=\frac{e_1(\kappa_1)\langle\Phi,\vec{a}\rangle+(\kappa_1+c)\gamma}{\alpha}+\gamma \kappa_1,\\
 0&=\frac{e_2(\kappa_1)\langle\Phi,\vec{a}+(\kappa_1+c)\mu}{\alpha}+\mu c.
 \end{split}
 \end{equation*}
 We obtain the expressions the functions $\gamma$ and $\mu$, namely,  
 $$\gamma=-\frac{e_1(\kappa_1)\langle\Phi,\vec{a}\rangle}{c+(1+\alpha)\kappa_1},\quad \mu=-\frac{e_2(\kappa_1)\langle\Phi,\vec{a}\rangle}{\kappa_1+(1+\alpha)c}.$$
  Differentiating $\mu$ with respect to $e_2$, we have
\begin{equation*}
e_2(\mu)=-\frac{\langle\Phi,\vec{a}\rangle}{\kappa_1+(1+\alpha)c}e_{22}(\kappa_1)+  \frac{2\langle\Phi,\vec{a}\rangle}{((1+\alpha)c+\kappa_1)^2}e_2(\kappa_1)^2.
\end{equation*}
On the other hand, the fourth equation of \eqref{d3} gives
$$e_2(\mu)=c\langle\Phi,\vec{a}\rangle \frac{c+\kappa_1}{\alpha},$$
because $\omega(e_2)=0$. From the last two expressions of $e_2(\mu)$ we obtain 
\begin{equation}\label{E1}
e_{22}(\kappa_1)=((\alpha+1)c+\kappa_1) \left(\frac{2 e_2(\kappa_1)^2}{((\alpha+1)c+\kappa_1)^2}-\frac{c (c+\kappa_1)}{\alpha }\right).
\end{equation}
 Inserting in \eqref{f1}, we have 
 $$\frac{(c-\kappa_1) \left(\alpha  \left(c^2+\kappa_1^2\right)+(c+\kappa_1)^2\right)}{\alpha }-\frac{2 (\alpha +2) e_2(\kappa_1)^2}{(\alpha+1)c+\kappa_1}=0.$$
 If $\alpha=-2$, then this equation is $(c+\kappa_1)^2-2 \left(c^2+\kappa_1^2\right)=0$, which implies that $\kappa_1$ is constant. If $\alpha\not=-2$, then 
\begin{equation}\label{E2}
e_2(\kappa_1)^2=\frac{(c-\kappa_1) ((\alpha+1)c+\kappa_1) \left((\alpha +1) c^2+2 c \kappa_1+(\alpha +1)\kappa_1^2\right)}{2 \alpha  (\alpha +2)}.
\end{equation}
If $e_2(\kappa_1)=0$ identically, then \eqref{E2} implies $(\alpha+1)c+\kappa_1=0$ or $(\alpha +1) c^2+2 c \kappa_1+(\alpha +1)\kappa_1^2=0$, hence $\kappa_1$ is constant. In both cases, we deduce that $\kappa_1$ is a constant, which it is a contradiction. Thus $e_2(\kappa_1)\not=0$.  Differentiate  with respect to $e_2$ in \eqref{E2}, and   simplify later  by $e_2(\kappa_1)$, we have 
 $$2 e_{22}(\kappa_1)= \frac{\left(-\alpha ^2+\alpha +2\right) c^3+2 (\alpha -1) \alpha  c^2 \kappa_1-3 \left(\alpha ^2+\alpha +2\right) c \kappa_1^2-4 (\alpha +1) \kappa_1^3}{2 \alpha  (\alpha +2)}.$$
Using the value of $e_{22}(\kappa_1)$ and $e_2(\kappa_1)^2$ in \eqref{E1} and \eqref{E2}, respectively, we have 
 $$-(\alpha -1) \kappa_1^2+2 (\alpha +1) c \kappa_1+(\alpha +1) c^2=0.$$
 This proves that $\kappa_1$ is constant, which it is a contradiction.
%%%%%%%%
\section{Proof of Theorem \ref{t3}}\label{s6}
%%%%

Since planes parallel to $\vec{a}$ and spheres centered at $\Pi$ are singular minimal surfaces with constant mean curvature, it remain to prove the converse. 

  Let $\Sigma$ be an $\alpha$-singular minimal surface with  constant mean curvature $H$.     Instead to use the previous computations of Sect. \ref{s2}, we do a direct proof. If a principal curvature is constant, then the other one is also constant because $H$ is constant. In such a case, the surface is isoparametric, hence   $\Sigma$ is a plane or a sphere. This proves the theorem in this situation. 
  
  From now, we suppose that both principal curvatures are not constant.   Let  $\Omega\subset\Sigma$ be an open set of $\Sigma$ formed by points $p$ where $ N(p)\not=\pm \vec{a}$. This open exists because $\Sigma$ is not a plane. We also assume that $\Omega$ is not formed by umbilic points because $\Sigma$ is not a plane neither a sphere. 
    The singular minimal surface equation \eqref{eq1} writes as 
$$H\langle\Phi(p),\vec{a}\rangle-\alpha\langle N(p),\vec{a}\rangle=0,\quad p\in \Omega.$$
Let $\{e_1,e_2\}$ be an orthonormal frame in $\Omega$ formed by principal directions, with $Se_i=\kappa_i e_i$. Differentiating the above identity  with respect to $e_1$ and $e_2$, we obtain 
$$(H+\alpha\kappa_1)\langle e_1,\vec{a}\rangle=0$$
$$(H+\alpha\kappa_2)\langle e_2,\vec{a}\rangle=0.$$
Since $N\not=\pm\vec{a}$ in $\Omega$, then $\langle e_1,\vec{a}\rangle\not=0$ or $\langle e_2,\vec{a}\rangle\not=0$ at some point   $p_0\in \Omega$. Without loss of generality, we suppose that $\langle e_1,\vec{a}\rangle\not=0$ at $p_0$. Then in an open set $\widetilde{\Omega}\subset\Omega$ around $p_0$ we have   $H+\alpha\kappa_1=0$. This proves that $\kappa_1$  is constant, which it is a contradiction.    
 

\subsection*{Acknowledgment}
The author  has been partially supported by MINECO/MICINN/FEDER grant no. PID2023-150727NB-I00,  and by the ``Mar\'{\i}a de Maeztu'' Excellence Unit IMAG, reference CEX2020-001105- M, funded by MCINN/AEI/10.13039/ 501100011033/ CEX2020-001105-M.
%%%%%%%%%%%%%%%%%%%%%%%%%%%%%%%%%%%%%%%%%%%%%%%%%%%%%%%%%%%%%%
\begin{thebibliography}{99}
  
  \bibitem{ay}   Aydin,  M. E.,   Erdur Kara, A.:  Ruled singular minimal surfaces. J. Geom. Phys. 195, Paper No. 105055, 8 pp. (2024)

\bibitem{bd}  Bemelmans, J., Dierkes, U.:  On a singular variational integral with linear growth, I: Existence and regularity of minimizers.   Arch. Ration. Mech. Anal.    100,  83--103  (1987) 


 
\bibitem{bht}   B\"{o}hme, R.,  Hildebrandt, S.,  Taush, E.:
The two-dimensional analogue of the catenary.
 Pacific J. Math. 88, 247--278 (1980)
 

\bibitem{di1}   Dierkes, U.:   A geometric maximum principle, Plateau's problem for surfaces of
prescribed mean curvature, and the two-dimensional analogue of the catenary. 
In: Hildebrandt, S.,   Leis, R., eds.  Partial Differential Equations and Calculus of Variations.  
Springer Lecture Notes in Mathematics, Vol. 1357,  pp. 116--141 (1988)

 
\bibitem{di2}   Dierkes, U.:  Singular minimal surfaces. In: Hildebrandt, S., Karcher, H., eds.  Geometric Analysis and Nonlinear Partial Differential Equations.  Berlin: Springer,   177--193 (2003)

\bibitem{di3} Dierkes, U.: On solutions of the singular minimal surface equation. Annali di Matematica 198(2),
505--516 (2019)

\bibitem{dg} Dierkes, U., Groh, N.: 
Symmetric solutions of the singular minimal surface equation.  
Ann. Global Anal. Geom. 60,  431--453  (2021)

\bibitem{dh}   Dierkes, U.,   Huisken, G.: The $n$-dimensional analogue of the catenary: existence and nonexistence.
 Pacific J. Math.      141, 47--54  (1990)

\bibitem{dl}   Dierkes, U.,   L\'opez, R.:  Cylindrical singular minimal surfaces. Atti Accad. Naz. Lincei Cl. Sci. Fis. Mat. Natur. 34, 547--576 (2023)

\bibitem{hn}   Hartman, P.    Nirenberg, L.: On spherical image maps whose Jacobians do not change sign. Amer. J. Math. 81, 901--920 (1959)
\bibitem{lo1}   L\'opez, R.:  Invariant singular minimal surfaces.  Ann. Global Anal. Geom.      53, 521--541 (2018)

\bibitem{lo3}  L\'opez, R.:  The Dirichlet problem for the $\alpha$-singular minimal surface equation.  Arch. Math. (Basel)   112, 213--222 (2019)

\bibitem{lo4}  L\'opez, R.:  Uniqueness of critical points and maximum principles of the singular minimal surface equation.  J. Differential Equations    266, 3927--3941  (2019)

\bibitem{lo2}   L\'opez, R.: Compact singular minimal surfaces with boundary.  Amer. J. Math.     142,  1771--1795 (2020)
   

\bibitem{ni}  Nitsche, J. C. C.:  A nonexistence theorem for the two-dimensional analogue of the catenary.  Analysis  6, 143--156 (1986)
 
 
 

\bibitem{wo} Wolfram Research, Inc., Mathematica, Version 10.0.0, Champaign, IL (2014)

 \end{thebibliography}

\end{document}



\documentclass[a4paper,12pt]{amsart}

  \usepackage{amssymb,amsthm}
 \setlength{\textwidth}{15cm}
\setlength{\oddsidemargin}{1cm}
\setlength{\evensidemargin}{1cm}
\setlength{\textheight}{21cm}
\setlength{\parskip}{2mm}
\setlength{\parindent}{0em}
\setlength{\headsep}{1.5cm}
%
%\usepackage{makeidx}   % allows index generation
\usepackage{graphicx,color}    % standard LaTeX graphics tool
  % when including figure files
%\usepackage{multicol}    % used for the two-column index

\usepackage{url} % <-- Para p\'aginas web o similar: \url{...}

\newtheorem{theorem}{Theorem}[section]
\newtheorem{corollary}[theorem]{Corollary}
\newtheorem{proposition}[theorem]{Proposition}
\newtheorem{lemma}[theorem]{Lemma}

\theoremstyle{definition}
\newtheorem{definition}[theorem]{Definition}
\newtheorem{example}[theorem]{Example}
\newtheorem{remark}[theorem]{Remark}


 \usepackage{showkeys}

\def\r{\mathbb R}
\def\h{\mathbb H}
\def\v{\vec{v}}
\def\e{\mathbf e}
\def\R{\mathbb R}
\def\n{\textbf{n}}
\def\a{\textbf{a}}
\def\s{\mathbb S}
\def\M{\mathbb M}
\begin{document}




\title{Singular minimal surfaces with constant  curvature}
\author{Rafael L\'opez}
\address{ Department of Geometry and Topology. University of Granada. 18071  Granada, Spain}
\email{rcamino@ugr.es}



\begin{abstract} 
We prove that spheres  are the only singular minimal surfaces with non-zero constant Gauss curvature. We also study singular minimal surfaces with zero constant Gauss curvature, proving that they must be cylindrical surfaces. We also classify singular minimal surfaces with one constant principal curvature and with constant mean curvature.    
\end{abstract}

